\section{Methods}
\label{sec:methods}

\subsection{Datasets}

The inorganic-crystal benchmark trains DPA4 on MPtrj, the Materials Project
trajectory corpus introduced with CHGNet. MPtrj contains relaxation and static
calculations for inorganic crystals across 89 elements, with energies, forces,
stresses and magnetic moments computed at the GGA or GGA+$U$
level~\cite{deng2023chgnet,jain2013materials}. Generalisation to metastability
ranking is evaluated on the Matbench Discovery benchmark and its WBM test
set~\cite{riebesell2025matbench,WBM}. The architectural ablations use the same
WBM-subsampled protocol as the main inorganic-crystal experiments so that
changes in accuracy and throughput reflect the controlled model component.

The molecular benchmark uses SPICE-MACE-OFF, the organic-molecule data set used
for MACE-OFF. It combines SPICE, molecular-dynamics conformations,
QMugs-derived larger molecules and water clusters; reference energies and
forces are evaluated at the $\omega$B97M-D3(BJ)/def2-TZVPPD level with
PSI4~\cite{kovacs2023mace,
    eastman2023spice,eastman2024nutmeg,isert2022qmugs,najibi2018nonlocal,
    grimme2010consistent,grimme2011effect,weigend2005balanced,rappoport2010property,
    smith2020psi4}.

\subsection{DPA4 model architecture}

Consider a system of $N$ atoms with atomic numbers $\{Z_i\}$ and positions
$R=\{\mathbf{R}_i\}_{i=1}^{N}$. DPA4 represents the potential energy as the sum
of a learned equivariant message-passing branch and an analytical
Ziegler--Biersack--Littmark (ZBL) short-range branch:
\begin{equation}
    E(Z,R)=E_{\Theta}^{\mathrm{NN}}\bigl(Z,R;\{\widetilde{\mathbf{r}}_{ij}\}\bigr)
    + E^{\mathrm{ZBL}}\bigl(Z,R;\{r_{ij}\}\bigr),
    \label{eq:dpa4-total-energy}
\end{equation}
where $r_{ij}=|\mathbf{R}_j-\mathbf{R}_i|$ is the true interatomic distance and
$\widetilde{\mathbf{r}}_{ij}$ is the clamped displacement consumed by the
learned branch inside the Native ZBL Zone Bridging window. Forces and virials are
obtained by differentiating the scalar energy in Eq.~\eqref{eq:dpa4-total-energy}
with respect to atomic positions and the cell, so the neural and analytical
branches define one conservative potential.

DPA4 carries node features in the real irreducible representation space
$V_{\leq L}\otimes\mathbb{R}^{C}$, where $V_{\leq L}=\bigoplus_{l=0}^{L}V_l$
and $V_l$ is the $(2l+1)$-dimensional irreducible representation of SO(3). The
$l=0$ subspace is invariant and is used for scalar readout; $l>0$ subspaces
transform equivariantly and carry angular information. Each interaction block
consists of a multi-focus SO(2) convolution in an edge-local gauge and an
equivariant feed-forward network. The architecture combines an Environment
Initial Embedding that injects a scalar local-environment prior at layer 0,
edge-conditioned radial degree coupling inside the local SO(2) frame,
envelope-gated softmax attention for smooth aggregation, a Lebedev $S^2$
activation in the feed-forward branch and Native ZBL Zone Bridging for
short-range repulsion. Geometry-dependent quantities, including radial basis
functions, cutoff envelopes, local gauges and bridging gates, are constructed
once per neighbour graph and reused across blocks.

\subsubsection{Initial equivariant representation}
\label{sec:initial-representation}

The initial scalar state is a learnable element embedding. For $l\geq1$, DPA4
adds a geometric initial embedding by projecting each neighbour direction into
real spherical harmonics and modulating it by radial-species features:
\begin{equation}
    \mathbf{h}_{i,\iota(l,m),c}^{(0)}
    \mathrel{+}=
    n_i\sum_{j:r_{ij}<r_{\mathrm{c}}}
    \eta_j\,Y_l^m(\widehat{\mathbf{r}}_{ij})\,
    \rho_{ij,l,c},\qquad l\geq 1.
    \label{eq:dpa4-gie}
\end{equation}
Here $\iota(l,m)$ indexes the coefficient of degree $l$ and order $m$,
$n_i$ is a smooth degree normalisation, $\eta_j$ is the source-freeze gate
defined below, and $\rho_{ij,l,c}$ is a smooth radial-species profile. This
initialisation is SO(3)-equivariant because it multiplies each degree-$l$
spherical harmonic by a scalar invariant.

DPA4 further conditions the $l=0$ slice with an Environment Initial Embedding.
A compact invariant descriptor $\mathcal{D}_i$ is built from a smooth local
environment matrix containing scalar and directional radial-species moments
(Supplementary Section~\ref{sec:sm-initial-embeddings}). This descriptor
provides an initialisation prior for the scalar backbone through feature-wise
affine modulation,
\begin{equation}
    \mathbf{h}_{i,l=0}^{(0)}
    \leftarrow
    \gamma(\mathcal{D}_i)\odot \mathbf{h}_{i,l=0}^{(0)}
    +\beta(\mathcal{D}_i).
    \label{eq:dpa4-film}
\end{equation}
The affine map is initialised near the identity, so the species embedding remains
the dominant signal at the beginning of training while the model learns when to
use environment conditioning.

\subsubsection{Local-frame SO(2) convolution}
\label{sec:so2}

Full SO(3)-equivariant tensor products require Clebsch--Gordan expansions whose
cost grows steeply with angular
order~\cite{batatia2022mace,liao2023equiformerv2}. DPA4 follows the local-frame
strategy of reducing SO(3) convolutions to SO(2)
operations~\cite{passaro2023reducing}: for each directed edge $(i,j)$, an
edge-local gauge $R_{ij}$ maps the bond direction to the reference axis,
\begin{equation}
    R_{ij}\widehat{\mathbf{r}}_{ij}=(0,0,1)^{\top}.
\end{equation}
The residual symmetry in this frame is the abelian group SO(2), so angular
orders $m$ can be processed in independent strata. DPA4 retains coefficients
with $|m|\leq M\leq L$ inside the edge-local convolution and applies a
degree-dependent lift factor after rotating back to compensate the norm loss
from truncation.

For a source feature $\mathbf{h}_j^{(b)}$ at block $b$, the edge-local tensor
is
\begin{equation}
    \mathbf{x}_{ij}=
    P_M D(R_{ij})\,\mathbf{h}_j^{(b)},
    \label{eq:dpa4-local-projection}
\end{equation}
where $D(R_{ij})$ is block diagonal in $l$ and $P_M$ selects the retained
$m$-strata. Edge-conditioned radial degree coupling then applies a smooth kernel
across angular degrees at fixed $|m|$:
\begin{equation}
    \mathbf{z}_{ij,l_{\mathrm{out}},m,c}
    =
    \sum_{l_{\mathrm{in}}\geq |m|}
    \mathcal{K}_{ij,l_{\mathrm{out}},l_{\mathrm{in}},|m|,c}\,
    \mathbf{x}_{ij,l_{\mathrm{in}},m,c}.
    \label{eq:dpa4-radial-coupling}
\end{equation}
In the DPA4 configurations using radial degree coupling, the kernel is a low-rank per-channel
function of scalar radial-species information and therefore does not introduce a
new equivariant edge feature. It generalises diagonal radial scaling while
preserving the local SO(2) selection rule: coefficients of different $|m|$ are
never mixed, and the $(-m,+m)$ pair transforms as one two-dimensional real
representation.

The subsequent SO(2)-equivariant linear maps are unrestricted on the $m=0$
subspace and use the real form of complex multiplication on each $|m|>0$
subspace. Multiple focus streams may process the same edge context in parallel;
a scalar-keyed softmax over the invariant $l=0$ slice assigns smooth
competition weights to these streams before the edge message is lifted back to
the global frame. The resulting message is
\begin{equation}
    \mathbf{m}_{ij}=
    \Xi_M D(R_{ij})^{\top}P_M^{\top}\,
    \mathcal{S}_{\Theta}(\mathbf{z}_{ij}),
    \label{eq:dpa4-edge-message}
\end{equation}
where $\mathcal{S}_{\Theta}$ denotes the multi-layer SO(2) stack and
$\Xi_M$ is the degree-wise truncation compensation.

\subsubsection{Envelope-gated softmax attention}
\label{sec:attn}

Neighbour aggregation is performed either by a smooth envelope-weighted sum or
by envelope-gated softmax attention. Attention logits are scalar functions of
the invariant $l=0$ destination and source features, plus a radial bias. For
focus $f$ and head $a$, let $\ell_{ij}^{(f,a)}$ be this invariant logit and
$\mu_i^{(f,a)}=\max_{k}\ell_{ik}^{(f,a)}$ over active neighbours. DPA4 uses
\begin{equation}
    \alpha_{ij}^{(f,a)}
    =
    \frac{s(r_{ij})^2\eta_j
        \exp\!\left(\ell_{ij}^{(f,a)}-\mu_i^{(f,a)}\right)}
    {\operatorname{softplus}(\zeta_{f,a})
        \exp\!\left(-\mu_i^{(f,a)}\right)
        +\sum_{k:r_{ik}<r_{\mathrm{c}}}
        s(r_{ik})^2\eta_k
        \exp\!\left(\ell_{ik}^{(f,a)}-\mu_i^{(f,a)}\right)}.
    \label{eq:dpa4-attention}
\end{equation}
The factor $s(r)^2$ sends the weight smoothly to zero at the cutoff, while the
positive denominator term prevents $0/0$ limits when all incident edges are
silenced by cutoff or bridging gates. Since all attention weights are invariant
scalars, multiplying them into $\mathbf{m}_{ij}$ preserves equivariance of every
higher-degree channel. A destination-side scalar gate further modulates each
attention head after aggregation.

\subsubsection{Equivariant feed-forward update and Lebedev \texorpdfstring{$S^2$}{S2} activation}
\label{sec:s2-method}

After each SO(2) convolution, DPA4 applies an equivariant feed-forward network
with residual connection. Degree-wise linear maps act independently on each
$l$-block, and higher-degree channels are gated by scalar functions of the
$l=0$ slice. In the reported DPA4 configurations, the feed-forward nonlinearity
uses a spherical-grid SwiGLU activation in the feed-forward branch. The
activation projects the $l\geq1$ value and gate tensors to quadrature points on
$S^2$, applies the nonlinearity pointwise on the sphere, and projects the
result back to irreducible coefficients; the scalar $l=0$ slice is processed by
the usual SwiGLU. DPA4 evaluates these spherical maps with Lebedev quadrature
rather than the product grids commonly used in Equiformer-family
implementations, reducing the number of quadrature points and the measured
equivariance error for the degrees used in this work (Supplementary
Section~\ref{sec:sm-s2}).

\subsubsection{Native ZBL Zone Bridging}
\label{sec:zbl_method}

The analytical branch uses the ZBL screened Coulomb pair
potential~\cite{ziegler1985},
\begin{equation}
    E_{ij}^{\mathrm{ZBL}}(r)=
    \frac{k_{\mathrm{e}}Z_iZ_j}{r}\,
    \Phi\!\left(\frac{r}{a_{ij}}\right),
    \qquad
    a_{ij}=\frac{0.88534\,a_0}{Z_i^{0.23}+Z_j^{0.23}},
    \label{eq:dpa4-zbl}
\end{equation}
with the universal screening function $\Phi$ given in Supplementary
Section~\ref{sec:sm-zone-bridging}. DPA4 does not treat this term as a
post-hoc wrapper. Instead, the learned branch and analytical branch are coupled
through Native ZBL Zone Bridging, which removes the direct learned short-range
pair channel while assigning the corresponding inner-zone pair interaction to
the analytical branch.

First, a $C^3$ clamped distance map $\widetilde{r}(r)$ connects a constant
inner zone to the identity outside the bridging window:
\begin{equation}
    \widetilde{r}(r)=
    \begin{cases}
        r_{\mathrm{in}}, & r\leq r_{\mathrm{in}},              \\
        r_{\mathrm{in}}+(r_{\mathrm{out}}-r_{\mathrm{in}})h_{\mathrm{c}}(t),
                         & r_{\mathrm{in}}<r<r_{\mathrm{out}}, \\
        r,               & r\geq r_{\mathrm{out}},
    \end{cases}
    \qquad
    t=\frac{r-r_{\mathrm{in}}}{r_{\mathrm{out}}-r_{\mathrm{in}}}.
    \label{eq:dpa4-clamp}
\end{equation}
The associated clamped displacement is
$\widetilde{\mathbf{r}}_{ij}=\widetilde{r}(r_{ij})\widehat{\mathbf{r}}_{ij}$.
The polynomial $h_{\mathrm{c}}$ is chosen to match a constant function at
$r_{\mathrm{in}}$ and the identity map at $r_{\mathrm{out}}$, including first
through third derivatives required for $C^3$ gluing.

Second, a source-freeze gate
\begin{equation}
    \eta_j=\prod_{i:r_{ji}<r_{\mathrm{c}}} w(r_{ji})
    \label{eq:dpa4-source-freeze}
\end{equation}
uses a separate $C^3$ switching amplitude $w$ that is exactly zero in the inner
zone and exactly one outside the bridging window. If atom $j$ has any neighbour
inside the frozen zone, then $\eta_j=0$ and all information sourced from $j$ is
removed from the geometric initial embedding, Environment Initial Embedding,
plain aggregation and attention aggregation. The clamp closes the scalar
distance channel, while the gate closes the direction and multi-hop propagation
channels. Thus the direct learned contribution associated with an inner-zone
pair has zero derivative with respect to that pair channel, while learned
many-body contributions from non-frozen neighbours remain part of the neural
potential.

\subsubsection{Symmetry guarantees}
\label{sec:symmetry}

DPA4 is invariant to translations because all geometric inputs are built from
relative displacements. It is permutation-equivariant at the feature level and
permutation-invariant at the energy level because all edge and node maps are
shared across atoms of the same species and neighbour aggregation is by
summation or scalar softmax aggregation. Rotational invariance follows because
the local-frame convolution, equivariant feed-forward network and initial
embedding all intertwine the SO(3) action on $V_{\leq L}$; the readout acts
only on the final $l=0$ invariant slice. Smoothness follows from the $C^3$
cutoff envelope, the $C^3$ bridging maps and the positive denominator in
Eq.~\eqref{eq:dpa4-attention}. Conservativeness follows from differentiating
the single scalar energy in Eq.~\eqref{eq:dpa4-total-energy}.

\subsection{Training}

DPA4 is trained by minimising a weighted sum of mean absolute errors in energy
and virials, together with the mean Euclidean norm of the atomic force-vector
residual:
\begin{equation}
    \mathcal{L}
    =
    \lambda_E
    \left\langle
    \frac{|E-\widehat{E}|}{N}
    \right\rangle
    +\lambda_F
    \left\langle
    \left\|\mathbf{F}_i-\widehat{\mathbf{F}}_i\right\|_2
    \right\rangle_i
    +\lambda_\Pi
    \left\langle
    \frac{\|\Pi-\widehat{\Pi}\|_1}{9N}
    \right\rangle .
    \label{eq:dpa4-loss}
\end{equation}
The benchmark models use the same training protocol: mixed-precision training with
full-precision geometric reductions, tensor-float matrix products where
available, a warmup--stable--decay (WSD) learning-rate schedule and the
HybridMuon optimiser~\cite{wen2024wsd,jordan2024muon,liu2025muon}. HybridMuon
routes matrix-valued hidden transformations to Muon updates and scalar,
normalisation or auxiliary parameters to Adam-family updates. Degree-indexed
equivariant tensors are updated in slice mode so that distinct representation
blocks are not flattened into a single matrix, and the Muon path uses
match-RMS scaling and Magma-lite alignment damping to keep the update scale
compatible with the Adam-family path~\cite{joo2026magma}. Complete model- and
dataset-specific hyperparameters are reported in Supplementary
Tables~\ref{tab:sm_main_ablation_config}--\ref{tab:sm_matbench_config}.

\subsection{Compiled conservative-force training}

DPA4 supports compiled conservative-force training. In conventional MLIP
training, force matching differentiates through the conservative force
$\mathbf{F}=-\partial E_{\Theta}/\partial\mathbf{R}$ and therefore requires the
mixed second derivative
$\partial^2E_{\Theta}/\partial\mathbf{R}\,\partial\Theta$, because the training
gradient differentiates through the coordinate derivative. This nested-autograd
structure has prevented previous conservative MLIP training loops from being
compiled end to end. DPA4 resolves this by tracing the energy-to-force
derivative with \texttt{make\_fx} into an ordinary tensor graph and lowering
the resulting graph with PyTorch Inductor~\cite{ansel2024pytorch2}. The
compiled graph contains both the energy evaluation and its coordinate
derivative, while the outer backward pass supplies the parameter derivative
required by force-loss training.

The implementation keeps the neighbour representation shape-stable and
preserves inactive neighbours as exactly silent contributions, so compiled
training retains the same conservative energy-to-force relationship as the
uncompiled model. In ablations, the compiled mixed-precision training path
gives the observed 2--3$\times$ wall-clock speedup without a systematic
accuracy penalty (Table~\ref{tab:dpa4_compile_ablation}).

% =============================================================
